% AI-produced by GPT-6 Astra (OpenAI, ChatGPT Pro mode). No priority claim.
% Independently audited in notes/finite_predimension_stabilizer_audit.md.
\subsection{Alternating amalgams and finite supports}

\begin{lemma}[Proved: a finite alternating bridge]
\label{A:predimension-alternating-bridge}
In either category $A$ of
Proposition~\ref{A:predimension-quasi}, let $B,C\leq_d D$ and put
$F=B\cap C$. Empty subsets are permitted as formal amalgamation
bases. Suppose $d_0,d_1:D\to E$ are arrows agreeing on $F$.
There is an arrow $E\to Z$ and a finite sequence of arrows
\[
 a_0,a_1,\ldots,a_{2N}:D\longrightarrow Z,
\]
whose endpoints are the given $d_0,d_1$ followed by $E\to Z$, such that
\[
 a_{2j}|_B=a_{2j+1}|_B,
 \qquad
 a_{2j+1}|_C=a_{2j+2}|_C
 \quad(0\leq j<N).
\]
\end{lemma}
\begin{proof}
If $D$ is empty, all asserted equalities are automatic and one
may take $Z=E$. Assume henceforth that $D$ is nonempty.
We construct an actual finite quotient and then prove positivity,
closedness, and orderability. Choose $N\geq4$, to be increased
below. Take one copy of $E$ and $2N-1$ copies $D_i$ of $D$,
for $1\leq i\leq2N-1$. Regard the $E$-copy as carrying $D_0$
by $d_0$ and $D_{2N}$ by $d_1$. Identify corresponding points
of $B$ between $D_i,D_{i+1}$ when $i$ is even, and corresponding
points of $C$ when $i$ is odd. At the two ends use the indicated
maps into $E$. The index graph is the circle of $s=2N$ factors
\[
 E,D_1,\ldots,D_{2N-1}.
\]
Let $Z$ be this strict set quotient, with exactly the union of the
copied hyperedges.

\paragraph{The fibers and induced subobjects.}
Each point of $F$ has one occurrence in every factor, connected
around the circle; the two endpoint identifications agree by
hypothesis. Every other point of an interior copy belongs to at
most one of $B,C$, so it crosses at most one consecutive overlap.
Its visibility away from $E$ has at most two indices. A point of
$E$ can lie in both $d_0(B)$ and $d_1(C)$, joining its two
neighbors through $E$; its visibility then has the three cyclically
consecutive indices $\{s-1,0,1\}$. Further propagation is impossible:
a point in $B\setminus C$ cannot cross a $C$-overlap, and conversely.
A point mapped to $F$ at an endpoint was already that same
$F$-point, by injectivity of the endpoint maps.

Thus every non-$F$ fiber is a proper cyclic interval of at most
three indices, and every $F$-fiber is the full circle. Every factor
map is injective, and the consecutive overlaps are exactly the
prescribed $B$- or $C$-images. Since $s\geq8$, these proper
intervals are shorter than half the circle. If the three vertices
of a copied edge are visible in another factor, they are visible
along the unique shortest index path between the two factors,
unless all three are full-visibility points, in which case any
index path works. At each overlap inducedness of its base
propagates the edge to the next factor. The bases are induced
both in $D$ and in their endpoint images in $E$. Hence every
factor is induced, and an edge's visibility is exactly the
intersection of its vertex visibility intervals. All edge
identifications, too, are generated by the neighboring relations.

\paragraph{Predimension and closedness.}
Put $m=\max(|B|,|C|)$ and increase $N$ so that
\[
 s=2N>3\bigl(\delta(E)+2m\bigr).
 \tag{AB}\label{eq:predimension-alternating-bound}
\]
Cut any overlap to produce a linear amalgam $L$. Every factor
$E,D_i$ is $d$-closed in $L$, by repeated free amalgamation.
The end bases are closed in $E$ by transitivity; all the other
bases are closed in the adjacent $D$-copies. The quotient
$q:L\to Z$ restores at most $m$ vertex-pair identifications.
For $U\subseteq Z$ and $\widetilde U=q^{-1}(U)$, it follows that
\[
 \delta_Z(U)\geq\delta_L(\widetilde U)-2m.
 \tag{3}\label{eq:predimension-alternating-loss}
\]
As in \eqref{eq:predimension-cut-loss}, the exact difference is
$2d_v-3d_e$, with $d_v\leq m$ and $d_e\geq0$.

If the proper visibility intervals of the vertices in $U$ miss
an index $r$, cut the overlap $(r-1,r)$. Every such vertex then
has a unique lift, and each full-circle fiber becomes a connected
path. Induced edges also have unique lifts. Thus
$Z[U]\cong L[\widetilde U]$. This gives positivity for every
nonempty such $U$, and strict relative positivity whenever $U$
properly contains any of the factors.

If those proper intervals cover all $s$ indices, let $J$ be the
set of indices in the linear amalgam where at least one lifted
non-$F$ vertex first appears. All vertices first appearing at $j$
are visible only among $j,j+1,j+2$. Therefore $3|J|\geq s$.
At each such step a selected vertex lies outside the incoming
closed overlap. Submodularity and strict closedness give an
integral predimension contribution of at least one, exactly as in
\eqref{eq:predimension-new-stage}. All other contributions are
nonnegative; the first nonempty factor subset likewise has
predimension at least one. Consequently
\[
 \delta_L(\widetilde U)\geq\lceil s/3\rceil,
 \qquad
 \delta_Z(U)\geq\lceil s/3\rceil-2m>\delta(E)\geq\delta(D)>0.
\]
The weak inequality uses $D\leq_d E$ at either endpoint.
Both cases prove that $Z$ is positive-predimension and that
$E$ and every $D_i$ are $d$-closed in it.

\paragraph{An actual ordered cocone.}
Embed $E$ increasingly in $\mathbb R$ and view $d_0,d_1$ as
coordinate maps on $D$. The linear interpolation
\[
 p_t(x)=(1-t)d_0(x)+td_1(x),\qquad 0\leq t\leq1,
\]
is strictly increasing for each $t$ and is constant on $F$.
Let $\Delta>0$ be a common lower bound for its adjacent-coordinate
gaps; the minimum of the two endpoint gap bounds suffices.
For $|D|\leq1$ choose any positive $\Delta$ instead. Put
$M_0=\max_x|d_1(x)-d_0(x)|$, and further increase $N$ until
$2M_0/N<\Delta$.

Set $b_{2j}=p_{j/N}$. Define $b_{2j+1}$ by keeping the
$B$-coordinates at $p_{j/N}$ and changing all coordinates outside
$B$ to $p_{(j+1)/N}$. This map differs coordinatewise from
$p_{(j+1)/N}$ by at most $M_0/N$, so each adjacent gap is at least
$\Delta-2M_0/N>0$. Hence it is increasing and agrees with
$b_{2j}$ on $B$. Now set $b_{2j+2}=p_{(j+1)/N}$. The intermediate
map already has the new coordinates on $C\setminus B$; those
on $B\cap C=F$ have not changed. Therefore it agrees with
$b_{2j+2}$ on $C$.

These maps, together with the coordinate map on $E$, satisfy
exactly the quotient's overlap equations, and have the required
endpoint values. They give a genuine function $Z\to\mathbb R$.
Order $Z$ by these values, breaking ties arbitrarily. Every factor
remains strictly increasing, since its coordinate map is strictly
increasing. This orders the already constructed hypergraph while
preserving the prescribed orders on all factors.

All lower bounds on $N$ are simultaneous lower bounds, so the
construction exists. In the point-and-edge-closed subcategory,
every atom of $Z$ lies in a closed old factor and is closed there;
transitivity puts $Z$ in the same subcategory. The factor maps are
the asserted sequence of arrows.
\end{proof}

\begin{proposition}[Proved: the stabilizer intersection formula]
\label{A:predimension-stabilizers}
Let $\mathbb M$ be a countable strong Fra\"iss\'e limit for either
category $A$ of Proposition~\ref{A:predimension-quasi}, and let
$G=\operatorname{Aut}(\mathbb M)$, preserving the order. For a
finite $d$-closed subset $B$, write $G_B$ for its pointwise
stabilizer. Then for all finite closed $B,C$,
\[
 \langle G_B,G_C\rangle=G_{B\cap C}.
\]
The equality is algebraic and in particular implies the corresponding
topological-generation assertion.
\end{proposition}
\begin{proof}
The category has countably many isomorphism types, joint embedding,
and strong amalgamation. The usual rich-sequence construction
therefore gives a countable strong limit; these are the hypotheses
and conclusion of \citet[Theorem~2.6]{EvansHubickaNesetril2019}.
Every finite subset of $\mathbb M$ lies in a finite closed
subobject, every isomorphism between finite closed subobjects extends
to an automorphism, and every finite closed extension of a finite
closed subobject embeds in $\mathbb M$ over that subobject.
For the last assertion, one may first freely amalgamate the proposed
extension with a sufficiently large member of the rich sequence
over the given subobject, and then use richness.

The generated subgroup plainly fixes $B\cap C$. Conversely let
$g\in G_{B\cap C}$. Choose a finite closed $D$ containing
$B\cup C$, and a finite closed $E$ containing $D\cup gD$.
Apply Lemma~\ref{A:predimension-alternating-bridge} to the two
embeddings $\operatorname{id}_D,g|_D:D\to E$. Embed its closed
extension $Z$ into $\mathbb M$ over $E$. Every intermediate
$D$-copy is closed, so its isomorphism from $D$ extends to an
automorphism $g_i$ of $\mathbb M$. Choose $g_0=1$ and
$g_{2N}=g$, which extend the stipulated endpoint copies. At each
step the two restrictions agree on $B$ or on $C$. Therefore
\[
 g_i^{-1}g_{i+1}\in
 \begin{cases}
 G_B,&i\text{ even},\\
 G_C,&i\text{ odd}.
 \end{cases}
\]
Multiplication telescopes to $g$, proving the reverse inclusion.
\end{proof}

\begin{corollary}[Proved: classification of open subgroups]
\label{A:predimension-open-subgroups}
Every open subgroup of $G$ is the pointwise stabilizer of a finite
$d$-closed subset of $\mathbb M$, possibly the empty subset.
Such a subset can be chosen contained in every finite closed
subset whose pointwise stabilizer is contained in the given group.
\end{corollary}
\begin{proof}
An open subgroup $U$ contains the pointwise stabilizer of a finite
set, hence that of a finite closed superset. Choose a finite closed
$B$ with $G_B\leq U$ of least cardinality. If $g\in U$, then
$G_{gB}=gG_Bg^{-1}\leq U$. Proposition~\ref{A:predimension-stabilizers}
gives $G_{B\cap gB}\leq U$. Minimality forces $gB=B$.
An increasing permutation of the finite ordered set $B$ is the
identity, so $g\in G_B$. This proves $U=G_B$.
If $G_C\leq U$ for another finite closed $C$, the same intersection
formula gives $G_{B\cap C}\leq U$, and minimality forces
$B\subseteq C$.
\end{proof}
